%! Least-Squares Regression — Unit 2, Lesson 2.4 — Student Packet Cover Sheet
\documentclass[12pt]{article}
\usepackage{apstats-article}
\usepackage{apstats-boxes}
\usepackage{ltablex}
\keepXColumns

\begin{document}

% Full-bleed navy banner. \coverbanner MEASURES the title block and sizes the
% band to it, so a lesson title long enough to wrap grows the band.
\coverbanner{Unit 2}{Lesson 2.4 \quad Least-Squares Regression}

\namedateperiod
\vspace{0.18in}

\boxguard[14]
\begin{learningtargetbox}
\small
% GRADUAL RELEASE: the vocabulary is taught in the guided notes, so the targets name it
% plainly. The AP Learning Objectives (DAT-1.G, DAT-1.H) stay in the teacher-facing plan.
\begin{enumerate}[leftmargin=1.5em, itemsep=5pt]
  \item \ldots{} explain what makes the \textbf{least-squares regression line} the best line,
        and compute its slope and intercept from $\bar x$, $s_x$, $\bar y$, $s_y$, and $r$.
  \item \ldots{} find the slope, the intercept, and $r^2$ in \textbf{computer regression
        output}.
  \item \ldots{} interpret the \textbf{slope} and the \textbf{$y$-intercept} in context ---
        and say when the intercept has no sensible meaning.
  \item \ldots{} interpret the \textbf{coefficient of determination} $r^2$ in context.
\end{enumerate}
\end{learningtargetbox}

\vspace{0.14in}

\boxguard[14]
\begin{tocbox}
\small
\renewcommand{\arraystretch}{1.5}
\begin{tabularx}{\linewidth}{@{} c l X r @{}}
\rowcolor{sky}
\textbf{\#} & \textbf{Component} & \textbf{Description} & \textbf{Score} \\
\midrule
1 & Warm-Up & A residual, two lines judged by their squared residuals, and the arithmetic
                 behind today's formula & \blank{1.2cm} \\
2 & Guided Notes \& Practice & Building the least-squares line for sleep and reaction time,
                 reading it off computer output, and used-car prices together & \blank{1.2cm} \\
% AP Practice is assigned but NOT scored: its score cell prints NA, never a \blank{}.
% Row 4 is the lesson's GRADED practice. Paper homework is the DEFAULT for this course; on the
% occasions DeltaMath is assigned instead, that replaces row 4 and its score cell is struck.
3 & AP Practice & Iced drinks and temperature: two multiple-choice items and one
                 free-response set --- practice, not a grade & \textbf{NA} \\
4 & Homework & Car weight and highway mileage: build the line, interpret it, and judge it
                 --- \textbf{scored, due the first class after two study halls} & \blank{1.2cm} \\
\midrule
  & \multicolumn{2}{r}{\textbf{Total}} & \blank{1.2cm} \\
\end{tabularx}
\end{tocbox}

\vspace{0.14in}

\boxguard[8]
\begin{remindbox}
\small
Of all the lines you could draw through a scatterplot, the \textbf{least-squares line} is the
one whose squared residuals add up to the least --- and it always runs through
$(\bar x, \bar y)$. Its slope describes the \textbf{predicted} value, not a promise about
every individual: real points scatter around the line, and a line fit to data we only
observed cannot show that $x$ \emph{causes} $y$. Before you read meaning into the intercept,
ask whether $x = 0$ makes sense at all.
\end{remindbox}

\end{document}
